\begin{proof}[Proof of \cref{obj:thm:zeng-fixed-positivity-polynomial-frontier}]
\begin{enumerate}
\item \textit{Uniform risk bound for the deterministic estimator.}
Let \(\overline C>0\) be the universal constant supplied by
\cref{obj:thm:unrestricted-mixed-count-upper}. Fix \(n,d\ge1\),
\(0<\epsilon<1/2\), and \(P_Z\in\mathcal D_{d,\epsilon}\).
Define the auxiliary assignment vector and synthetic sample by
\[
\mathbf u\sim\operatorname{Unif}(\{0,1\}^n),
\qquad
\mathbf O_n^{\mathrm{syn}}
:=(o_i(u_i))_{i=1}^n,
\]
with \(\mathbf u\) independent of \(\mathbf Z_n\).
By \cref{obj:lem:synthetic-randomization-kernel}, there exists
\(P_Z^*\in\mathcal M^+_{n,d,\epsilon}\) such that
\[
\mathcal L_{P_Z,\mathbf u}(\mathbf O_n^{\mathrm{syn}})
=
\bigl(\mathcal L_{P_Z^*}(X,A,S,R,RY)\bigr)^{\otimes n},
\qquad
\tau(P_Z^*)=\theta_Z(P_Z).
\]
Thus the synthetic sample satisfies the sampling and model hypotheses of
\cref{obj:thm:unrestricted-mixed-count-upper} at \(q=\epsilon\).

For each fixed sample realization
\(\mathbf z=((x_i,b_i,y_i))_{i=1}^n\), Cauchy--Schwarz on the
\(2^n\) assignment vectors gives
\[
\begin{split}
&\left[
\sum_{\mathbf u\in\{0,1\}^n}
\left\{
\widehat\tau^{\mathrm{mix}}_{n,d,\epsilon}
\bigl(o_1(u_1),\ldots,o_n(u_n)\bigr)
-\theta_Z(P_Z)
\right\}
\right]^2\\
&\qquad\le
2^n\sum_{\mathbf u\in\{0,1\}^n}
\left\{
\widehat\tau^{\mathrm{mix}}_{n,d,\epsilon}
\bigl(o_1(u_1),\ldots,o_n(u_n)\bigr)
-\theta_Z(P_Z)
\right\}^2.
\end{split}
\]
Dividing by \(2^{2n}\) and using the defining finite average yields
\[
\begin{split}
&\left(
\widehat\theta^{\mathrm{poly}}_{n,d,\epsilon}(\mathbf z)
-\theta_Z(P_Z)
\right)^2\\
&\qquad\le
\frac1{2^n}
\sum_{\mathbf u\in\{0,1\}^n}
\left\{
\widehat\tau^{\mathrm{mix}}_{n,d,\epsilon}
\bigl(o_1(u_1),\ldots,o_n(u_n)\bigr)
-\theta_Z(P_Z)
\right\}^2.
\end{split}
\]
Integrating this inequality, identifying the uniform coin average with
the synthetic sample law, and then applying
\cref{obj:thm:unrestricted-mixed-count-upper}, we obtain
\[
\begin{split}
&E_{P_Z}\!\left[
\left(
\widehat\theta^{\mathrm{poly}}_{n,d,\epsilon}(\mathbf Z_n)
-\theta_Z(P_Z)
\right)^2
\right]\\
&\quad\le
E_{P_Z,\mathbf u}\!\left[
\left(
\widehat\tau^{\mathrm{mix}}_{n,d,\epsilon}
(\mathbf O_n^{\mathrm{syn}})
-\theta_Z(P_Z)
\right)^2
\right]\\
&\quad=
E_{P_Z^*}\!\left[
\left(
\widehat\tau^{\mathrm{mix}}_{n,d,\epsilon}(\mathbf O_n)
-\tau(P_Z^*)
\right)^2
\right]
\le \mathfrak U_{n,d,\epsilon}
\le \overline C\,r_{n,d,\epsilon}.
\end{split}
\]
Taking the supremum over a nonempty observational class proves the
first assertion. If that class is empty, its real-valued supremum is
zero, and the same bound follows from
\[
0\le r_{n,d,\epsilon}
=
\min\left\{
1,\frac1{n\epsilon}
+\left[\frac d{n\epsilon\log(e+n\epsilon)}\right]^2
\right\},
\qquad
0<\overline C.
\]
% lean: CausalSmith.Stat.MarRareqLogfrontier.thetaPoly_uniform_upper

\item \textit{Upper bound at fixed positivity.}
Fix \(0<\epsilon<1/2\). For \(n,d\ge1\), define
\[
r^{\mathrm{raw}}_{n,d}
:=
\frac1n+\left[\frac d{n\log(e+n)}\right]^2,
\qquad
r^{\mathrm{fix}}_{n,d}:=\min\{1,r^{\mathrm{raw}}_{n,d}\},
\]
and set
\[
\ell_{\epsilon,n}:=\log(e+n\epsilon),
\qquad
\ell_{1,n}:=\log(e+n),
\qquad
\gamma_{\epsilon,\mathrm{cmp}}
:=\frac{1-\log\epsilon}{\epsilon}.
\]
All denominators below are positive. In particular,
\[
\ell_{\epsilon,n}\ge1,
\qquad
\ell_{1,n}>0,
\qquad
\log\epsilon\le0,
\qquad
\gamma_{\epsilon,\mathrm{cmp}}\ge1.
\]
Since \(\epsilon\le1\),
\[
e+n\le\frac{e+n\epsilon}{\epsilon}.
\]
Monotonicity of the logarithm therefore gives
\[
\ell_{1,n}
\le \ell_{\epsilon,n}-\log\epsilon
\le (1-\log\epsilon)\ell_{\epsilon,n},
\]
where the second inequality follows from
\[
(-\log\epsilon)(\ell_{\epsilon,n}-1)\ge0.
\]
Consequently,
\[
n\ell_{1,n}
\le
\gamma_{\epsilon,\mathrm{cmp}}\,
n\epsilon\ell_{\epsilon,n}.
\]
Also,
\[
\epsilon\gamma_{\epsilon,\mathrm{cmp}}^2
\ge
\epsilon\gamma_{\epsilon,\mathrm{cmp}}
=
1-\log\epsilon
\ge1.
\]
These two comparisons imply, respectively,
\[
\left[\frac d{n\epsilon\ell_{\epsilon,n}}\right]^2
\le
\gamma_{\epsilon,\mathrm{cmp}}^2
\left[\frac d{n\ell_{1,n}}\right]^2,
\qquad
\frac1{n\epsilon}
\le
\gamma_{\epsilon,\mathrm{cmp}}^2\frac1n.
\]
Adding gives
\[
\frac1{n\epsilon}
+\left[\frac d{n\epsilon\ell_{\epsilon,n}}\right]^2
\le
\gamma_{\epsilon,\mathrm{cmp}}^2 r^{\mathrm{raw}}_{n,d}.
\]
Truncation preserves the needed comparison. Indeed, if
\(r^{\mathrm{raw}}_{n,d}\le1\), then
\[
\min\{1,\gamma_{\epsilon,\mathrm{cmp}}^2r^{\mathrm{raw}}_{n,d}\}
\le
\gamma_{\epsilon,\mathrm{cmp}}^2r^{\mathrm{raw}}_{n,d}
=
\gamma_{\epsilon,\mathrm{cmp}}^2r^{\mathrm{fix}}_{n,d};
\]
if \(r^{\mathrm{raw}}_{n,d}\ge1\), then
\[
\min\{1,\gamma_{\epsilon,\mathrm{cmp}}^2r^{\mathrm{raw}}_{n,d}\}
\le1
\le
\gamma_{\epsilon,\mathrm{cmp}}^2
=
\gamma_{\epsilon,\mathrm{cmp}}^2r^{\mathrm{fix}}_{n,d}.
\]
Thus
\[
r_{n,d,\epsilon}
\le
\gamma_{\epsilon,\mathrm{cmp}}^2r^{\mathrm{fix}}_{n,d}.
\]

The mixed-count estimator takes values in \([-1,1]\), by its clipped
output and averaging construction in \cref{obj:def:mixed-count-estimator}.
Hence, for every sample,
\[
-2^n
\le
\sum_{\mathbf u\in\{0,1\}^n}
\widehat\tau^{\mathrm{mix}}_{n,d,\epsilon}
\bigl(o_1(u_1),\ldots,o_n(u_n)\bigr)
\le2^n,
\]
so
\[
-1\le
\widehat\theta^{\mathrm{poly}}_{n,d,\epsilon}(\mathbf z)
\le1.
\]
This finite-sample map is measurable and is an admissible deterministic
procedure in \cref{obj:def:zeng-minimax-risk}. The defining infimum and
the preceding bounds therefore give
\[
\begin{split}
\mathfrak R^Z_{n,d,\epsilon}
&\le
\sup_{P_Z\in\mathcal D_{d,\epsilon}}
E_{P_Z}\!\left[
\left(
\widehat\theta^{\mathrm{poly}}_{n,d,\epsilon}(\mathbf Z_n)
-\theta_Z(P_Z)
\right)^2
\right]\\
&\le \overline C\,r_{n,d,\epsilon}
\le C_{\epsilon,\mathrm{up}}\,r^{\mathrm{fix}}_{n,d},
\qquad
C_{\epsilon,\mathrm{up}}
:=
\overline C\,\gamma_{\epsilon,\mathrm{cmp}}^2>0.
\end{split}
\]
% lean: CausalSmith.Stat.MarRareqLogfrontier.zeng_minimax_upper_fixedQ

\item \textit{Published lower bound on the zero-control subclass.}
Define
\[
\mathcal D^{(0)}_{d,\epsilon}
:=
\left\{
P_Z\in\mathcal D_{d,\epsilon}:
\mu_{0x}=0\ \text{for every }x\text{ with }p_x>0
\right\},
\]
and its minimax risk by
\[
\mathfrak R^{Z,0}_{n,d,\epsilon}
:=
\inf_{\widehat\theta}
\sup_{P_Z\in\mathcal D^{(0)}_{d,\epsilon}}
E_{P_Z}\!\left[
\bigl(\widehat\theta-\theta_Z(P_Z)\bigr)^2
\right],
\]
using the same procedure class as in
\cref{obj:def:zeng-minimax-risk}. The subclass inclusion gives
\[
\mathfrak R^{Z,0}_{n,d,\epsilon}
\le
\mathfrak R^Z_{n,d,\epsilon}.
\]
The published lower bound
\citep[Theorem~2, Section~3.2, p.~9, and its zero-control construction
in Section~C.5, pp.~39--44]{ZengBalakrishnanHanKennedy2026}
supplies constants
\[
c_{\epsilon,\mathrm{lit}}>0,
\qquad
c_{\epsilon,\mathrm{dim}}>0,
\qquad
n_{\epsilon,\mathrm{lit}}\in\mathbb N
\]
such that
\[
\mathfrak R^{Z,0}_{n,d,\epsilon}
\ge
c_{\epsilon,\mathrm{lit}}
\left\{
\frac1n+\left[\frac d{n\log n}\right]^2
\right\}
\]
whenever
\[
n\ge n_{\epsilon,\mathrm{lit}},
\qquad
1\le d\le c_{\epsilon,\mathrm{dim}}\,n\log n.
\]
The target belongs to \([-1,1]\), since it is a probability-weighted
average of differences of two means in \([0,1]\). By
\cref{obj:lem:randomized-kernel-barycenter}, clipping and then taking
conditional barycenters preserve or reduce squared-error risk at
every law. Thus the deterministic and randomized minimax values agree,
and the published bound applies to the procedure class above.
% lean: CausalSmith.Stat.MarRareqLogfrontier.ZengTheoremTwoLower

\item \textit{A uniform parametric lower bound.}
We next establish a universal constant \(c_{\mathrm{par}}>0\) such that
\[
\frac{c_{\mathrm{par}}}{n}
\le
\mathfrak R^{Z,0}_{n,d,\epsilon}
\le
\mathfrak R^Z_{n,d,\epsilon}
\qquad
(n,d\ge1,\ 0<\epsilon<1/2).
\]
For \(n\ge1\), define
\[
\delta_n:=\frac1{4\sqrt n},
\qquad
\vartheta^{\mathrm{test}}_0:=\frac12,
\qquad
\vartheta^{\mathrm{test}}_1:=\frac12+\delta_n.
\]
Then \(0<\delta_n\le1/4\), so both target values lie in \([0,1]\).
For \(\iota\in\{0,1\}\), define the one-record law
\(P^{\mathrm{test}}_{Z,\iota}\) on
\(\{1,\ldots,d\}\times\{0,1\}^2\) by
\[
P^{\mathrm{test}}_{Z,\iota}\{(x,b,y)\}
=
\begin{cases}
1/2,&(x,b,y)=(1,0,0),\\
(1-\vartheta^{\mathrm{test}}_\iota)/2,
 &(x,b,y)=(1,1,0),\\
\vartheta^{\mathrm{test}}_\iota/2,
 &(x,b,y)=(1,1,1),\\
0,&\text{otherwise}.
\end{cases}
\]
Its sole occupied baseline category has treatment probability \(1/2\),
control mean zero, and treated mean \(\vartheta^{\mathrm{test}}_\iota\).
Thus
\[
P^{\mathrm{test}}_{Z,\iota}\in\mathcal D^{(0)}_{d,\epsilon},
\qquad
\theta_Z(P^{\mathrm{test}}_{Z,\iota})
=\vartheta^{\mathrm{test}}_\iota.
\]

The two laws have common support. Define their chi-square divergence by
\[
\chi^2\!\left(
P^{\mathrm{test}}_{Z,1}\Vert P^{\mathrm{test}}_{Z,0}
\right)
:=
\sum_{\substack{z\in\{1,\ldots,d\}\times\{0,1\}^2:\\
P^{\mathrm{test}}_{Z,0}\{z\}>0}}
\frac{
\bigl(P^{\mathrm{test}}_{Z,1}\{z\}
-P^{\mathrm{test}}_{Z,0}\{z\}\bigr)^2
}{
P^{\mathrm{test}}_{Z,0}\{z\}
}.
\]
Only the two treated-outcome probabilities change. Each changes in
absolute value by \(\delta_n/2\), and each has reference probability
\(1/4\). Hence
\[
\chi^2\!\left(
P^{\mathrm{test}}_{Z,1}\Vert P^{\mathrm{test}}_{Z,0}
\right)
=2\delta_n^2
=\frac1{8n}.
\]
Write the sample laws and their likelihood ratio as
\[
Q^{\mathrm{test}}_\iota
:=
\bigl(P^{\mathrm{test}}_{Z,\iota}\bigr)^{\otimes n},
\qquad \iota\in\{0,1\},
\]
\[
L_{\mathrm{test}}(\mathbf z)
:=
\begin{cases}
Q^{\mathrm{test}}_1\{\mathbf z\}/
Q^{\mathrm{test}}_0\{\mathbf z\},
 &Q^{\mathrm{test}}_0\{\mathbf z\}>0,\\
0,&Q^{\mathrm{test}}_0\{\mathbf z\}=0.
\end{cases}
\]
Independence gives the second-moment identity and bound
\[
E_{Q^{\mathrm{test}}_0}[L_{\mathrm{test}}^2]
=
\left[
1+\chi^2\!\left(
P^{\mathrm{test}}_{Z,1}\Vert P^{\mathrm{test}}_{Z,0}
\right)
\right]^n
=
\left(1+\frac1{8n}\right)^n
\le e^{1/8}.
\]

Define the universal testing constant by
\[
c_{\mathrm{test}}:=\frac1{4e^{1/8}}>0.
\]
For any event \(\mathcal A\) in the finite sample space, set
\[
\eta_{\mathcal A}
:=
Q^{\mathrm{test}}_1(\mathcal A^c)
+
Q^{\mathrm{test}}_0(\mathcal A).
\]
Cauchy--Schwarz yields
\[
1-\eta_{\mathcal A}
\le Q^{\mathrm{test}}_1(\mathcal A)
=
E_{Q^{\mathrm{test}}_0}
[L_{\mathrm{test}}\mathbf1_{\mathcal A}]
\le
\sqrt{e^{1/8}Q^{\mathrm{test}}_0(\mathcal A)}.
\]
If \(\eta_{\mathcal A}<c_{\mathrm{test}}\), then
\(Q^{\mathrm{test}}_0(\mathcal A)\le\eta_{\mathcal A}\) implies
\[
\sqrt{e^{1/8}Q^{\mathrm{test}}_0(\mathcal A)}
<\frac12,
\qquad
1-\eta_{\mathcal A}>\frac34,
\]
because \(c_{\mathrm{test}}\le1/4\). This contradicts the preceding
inequality. Therefore every event satisfies
\[
Q^{\mathrm{test}}_1(\mathcal A^c)
+
Q^{\mathrm{test}}_0(\mathcal A)
\ge c_{\mathrm{test}}.
\]

Consider an arbitrary estimator. By
\cref{obj:lem:randomized-kernel-barycenter}, it suffices to consider its
clipped conditional action kernel \(\kappa\), supported on \([-1,1]\).
Define its barycenter by
\[
b_\kappa(\mathbf z)
:=
\int_{[-1,1]}h\,\kappa(\mathbf z,dh).
\]
The same cited result gives
\[
b_\kappa(\mathbf z)\in[-1,1],
\qquad
E_{Q^{\mathrm{test}}_\iota}
\left[
(b_\kappa-\vartheta^{\mathrm{test}}_\iota)^2
\right]
\le
E_{P^{\mathrm{test}}_{Z,\iota}}
\left[
(\widehat\theta-\vartheta^{\mathrm{test}}_\iota)^2
\right],
\qquad \iota\in\{0,1\},
\]
where the right-hand expectation includes the estimator's
randomization.

Define the nearest-target test and its large-error events by
\[
\mathcal A_{\mathrm{test}}
:=
\left\{
\mathbf z:
|b_\kappa(\mathbf z)-\vartheta^{\mathrm{test}}_1|
\le
|b_\kappa(\mathbf z)-\vartheta^{\mathrm{test}}_0|
\right\},
\]
\[
\mathcal A_{\mathrm{bad},\iota}
:=
\left\{
\mathbf z:
|b_\kappa(\mathbf z)-\vartheta^{\mathrm{test}}_\iota|
\ge\frac{\delta_n}{2}
\right\},
\qquad \iota\in\{0,1\}.
\]
The triangle inequality gives
\[
\delta_n
=
|\vartheta^{\mathrm{test}}_1-\vartheta^{\mathrm{test}}_0|
\le
|b_\kappa-\vartheta^{\mathrm{test}}_0|
+
|b_\kappa-\vartheta^{\mathrm{test}}_1|.
\]
Consequently,
\[
\mathcal A_{\mathrm{test}}\subseteq\mathcal A_{\mathrm{bad},0},
\qquad
\mathcal A_{\mathrm{test}}^c\subseteq\mathcal A_{\mathrm{bad},1}.
\]
Applying the testing bound and then these inclusions yields
\[
c_{\mathrm{test}}
\le
Q^{\mathrm{test}}_1(\mathcal A_{\mathrm{bad},1})
+
Q^{\mathrm{test}}_0(\mathcal A_{\mathrm{bad},0}),
\]
and hence
\[
\frac{c_{\mathrm{test}}}{2}
\le
\max_{\iota\in\{0,1\}}
Q^{\mathrm{test}}_\iota(\mathcal A_{\mathrm{bad},\iota}).
\]
On each large-error event, squared loss is at least
\((\delta_n/2)^2\). Therefore
\[
\begin{split}
\frac{c_{\mathrm{test}}\delta_n^2}{8}
&\le
\left(\frac{\delta_n}{2}\right)^2
\max_{\iota\in\{0,1\}}
Q^{\mathrm{test}}_\iota(\mathcal A_{\mathrm{bad},\iota})\\
&\le
\max_{\iota\in\{0,1\}}
E_{Q^{\mathrm{test}}_\iota}
\left[
(b_\kappa-\vartheta^{\mathrm{test}}_\iota)^2
\right]\\
&\le
\max_{\iota\in\{0,1\}}
E_{P^{\mathrm{test}}_{Z,\iota}}
\left[
(\widehat\theta-\vartheta^{\mathrm{test}}_\iota)^2
\right].
\end{split}
\]
Both testing laws belong to the zero-control subclass, so taking the
minimax infimum gives
\[
\frac{c_{\mathrm{par}}}{n}
=
\frac{c_{\mathrm{test}}\delta_n^2}{8}
\le
\mathfrak R^{Z,0}_{n,d,\epsilon}
\le
\mathfrak R^Z_{n,d,\epsilon},
\qquad
c_{\mathrm{par}}:=\frac{c_{\mathrm{test}}}{128}>0.
\]
% lean: CausalSmith.Stat.MarRareqLogfrontier.zeng_parametric_minimax_lower

\item \textit{Calibration and the large-sample bound within the published range.}
Choose an integer \(n_{\epsilon,1}\) satisfying
\[
n_{\epsilon,1}\ge\frac2{c_{\epsilon,\mathrm{dim}}},
\]
and define
\[
n_{\epsilon,*}
:=
\max\{n_{\epsilon,\mathrm{lit}},3,n_{\epsilon,1}\},
\qquad
c_{\epsilon,\mathrm{sat}}
:=
c_{\epsilon,\mathrm{lit}}
\left(\frac{c_{\epsilon,\mathrm{dim}}}{2}\right)^2,
\]
\[
c_\epsilon
:=
\min\left\{
\frac{c_{\mathrm{par}}}{n_{\epsilon,*}},
\min\{c_{\epsilon,\mathrm{lit}},c_{\epsilon,\mathrm{sat}}\}
\right\}>0.
\]
For all \(n,d\ge1\),
\[
0\le r^{\mathrm{fix}}_{n,d}\le1.
\]

First suppose \(n\ge n_{\epsilon,*}\). Then
\[
n\ge n_{\epsilon,\mathrm{lit}},
\qquad
n\ge3,
\qquad
\log n>1,
\qquad
n\log n>0.
\]
Furthermore, the choice of \(n_{\epsilon,1}\) gives
\[
2\le c_{\epsilon,\mathrm{dim}}\,n
\le c_{\epsilon,\mathrm{dim}}\,n\log n.
\]
If
\[
d\le c_{\epsilon,\mathrm{dim}}\,n\log n,
\]
the published bound and subclass inclusion imply
\[
c_{\epsilon,\mathrm{lit}}
\left\{
\frac1n+\left[\frac d{n\log n}\right]^2
\right\}
\le
\mathfrak R^{Z,0}_{n,d,\epsilon}
\le
\mathfrak R^Z_{n,d,\epsilon}.
\]
Since
\[
0<\log n\le\log(e+n),
\qquad
0\le\frac d{n\log(e+n)}
\le\frac d{n\log n},
\]
we have
\[
r^{\mathrm{fix}}_{n,d}
\le
\frac1n+\left[\frac d{n\log(e+n)}\right]^2
\le
\frac1n+\left[\frac d{n\log n}\right]^2.
\]
Using \(c_\epsilon\le c_{\epsilon,\mathrm{lit}}\) gives
\[
\begin{split}
c_\epsilon r^{\mathrm{fix}}_{n,d}
&\le
c_\epsilon
\left\{
\frac1n+\left[\frac d{n\log n}\right]^2
\right\}\\
&\le
c_{\epsilon,\mathrm{lit}}
\left\{
\frac1n+\left[\frac d{n\log n}\right]^2
\right\}
\le
\mathfrak R^Z_{n,d,\epsilon}.
\end{split}
\]
% lean: CausalSmith.Stat.MarRareqLogfrontier.zeng_minimax_lower_fixedQ

\item \textit{The large-alphabet case.}
Continue to assume \(n\ge n_{\epsilon,*}\), and now suppose
\[
d>c_{\epsilon,\mathrm{dim}}\,n\log n.
\]
Define the smaller alphabet size by
\[
d_{\epsilon,n}^{\mathrm{pad}}
:=
\left\lfloor c_{\epsilon,\mathrm{dim}}\,n\log n\right\rfloor.
\]
The preceding lower bound of two on the quantity inside the floor gives
\[
1\le d_{\epsilon,n}^{\mathrm{pad}}\le d,
\qquad
d_{\epsilon,n}^{\mathrm{pad}}
\le c_{\epsilon,\mathrm{dim}}\,n\log n.
\]
Also, the elementary floor bound gives
\[
\begin{split}
\frac{c_{\epsilon,\mathrm{dim}}}{2}\,n\log n
&\le c_{\epsilon,\mathrm{dim}}\,n\log n-1\\
&<d_{\epsilon,n}^{\mathrm{pad}}.
\end{split}
\]
The published lower bound therefore applies at this smaller alphabet:
\[
c_{\epsilon,\mathrm{lit}}
\left\{
\frac1n+
\left[
\frac{d_{\epsilon,n}^{\mathrm{pad}}}{n\log n}
\right]^2
\right\}
\le
\mathfrak R^{Z,0}_{n,d_{\epsilon,n}^{\mathrm{pad}},\epsilon}
\le
\mathfrak R^Z_{n,d_{\epsilon,n}^{\mathrm{pad}},\epsilon}.
\]

To compare the two alphabet sizes, extend each law
\(P_Z\in\mathcal D_{d_{\epsilon,n}^{\mathrm{pad}},\epsilon}\)
by zero masses. Explicitly, define the \(d\)-category law
\(\operatorname{pad}_{d_{\epsilon,n}^{\mathrm{pad}},d}P_Z\) by
\[
\bigl(\operatorname{pad}_{d_{\epsilon,n}^{\mathrm{pad}},d}P_Z\bigr)
\{(x,b,y)\}
:=
\begin{cases}
P_Z\{(x,b,y)\},&1\le x\le d_{\epsilon,n}^{\mathrm{pad}},\\
0,&d_{\epsilon,n}^{\mathrm{pad}}<x\le d,
\end{cases}
\qquad b,y\in\{0,1\}.
\]
On the original categories, all baseline, treatment, and outcome
probabilities are preserved; the added categories have zero mass.
Thus
\[
\operatorname{pad}_{d_{\epsilon,n}^{\mathrm{pad}},d}P_Z
\in\mathcal D_{d,\epsilon},
\qquad
\theta_Z\!\left(
\operatorname{pad}_{d_{\epsilon,n}^{\mathrm{pad}},d}P_Z
\right)
=\theta_Z(P_Z).
\]
For a \(d\)-category estimator \(\widehat\theta\), define
\(\widehat\theta^{\mathrm{res}}\) on the smaller sample space by
\[
\widehat\theta^{\mathrm{res}}
\bigl((x_i,b_i,y_i)_{i=1}^n\bigr)
:=
\widehat\theta
\bigl((x_i,b_i,y_i)_{i=1}^n\bigr),
\qquad
x_i\in\{1,\ldots,d_{\epsilon,n}^{\mathrm{pad}}\},
\]
preserving its conditional action distribution if it is randomized.
The product sample law is preserved under this inclusion, so
\[
\begin{split}
&E_{P_Z}\!\left[
\left(\widehat\theta^{\mathrm{res}}-\theta_Z(P_Z)\right)^2
\right]\\
&\qquad=
E_{\operatorname{pad}_{d_{\epsilon,n}^{\mathrm{pad}},d}P_Z}
\!\left[
\left(
\widehat\theta-
\theta_Z\!\left(
\operatorname{pad}_{d_{\epsilon,n}^{\mathrm{pad}},d}P_Z
\right)
\right)^2
\right].
\end{split}
\]
Taking suprema and then infima proves
\[
\mathfrak R^Z_{n,d_{\epsilon,n}^{\mathrm{pad}},\epsilon}
\le
\mathfrak R^Z_{n,d,\epsilon}.
\]

Since \(n\log n>0\), the floor bound also gives
\[
\frac{c_{\epsilon,\mathrm{dim}}}{2}
\le
\frac{d_{\epsilon,n}^{\mathrm{pad}}}{n\log n}.
\]
Squaring and using \(1/n\ge0\), we obtain
\[
c_{\epsilon,\mathrm{sat}}
\le
c_{\epsilon,\mathrm{lit}}
\left\{
\frac1n+
\left[
\frac{d_{\epsilon,n}^{\mathrm{pad}}}{n\log n}
\right]^2
\right\}.
\]
Finally, \(r^{\mathrm{fix}}_{n,d}\le1\) and
\(c_\epsilon\le c_{\epsilon,\mathrm{sat}}\) give
\[
\begin{split}
c_\epsilon r^{\mathrm{fix}}_{n,d}
&\le c_\epsilon
\le c_{\epsilon,\mathrm{sat}}\\
&\le
c_{\epsilon,\mathrm{lit}}
\left\{
\frac1n+
\left[
\frac{d_{\epsilon,n}^{\mathrm{pad}}}{n\log n}
\right]^2
\right\}\\
&\le
\mathfrak R^{Z,0}_{n,d_{\epsilon,n}^{\mathrm{pad}},\epsilon}
\le
\mathfrak R^Z_{n,d_{\epsilon,n}^{\mathrm{pad}},\epsilon}
\le
\mathfrak R^Z_{n,d,\epsilon}.
\end{split}
\]
% lean: CausalSmith.Stat.MarRareqLogfrontier.zeng_minimax_lower_fixedQ

\item \textit{Small samples and the final constants.}
If \(1\le n<n_{\epsilon,*}\), the parametric bound and the calibration
of \(c_\epsilon\) imply
\[
c_\epsilon r^{\mathrm{fix}}_{n,d}
\le c_\epsilon
\le\frac{c_{\mathrm{par}}}{n_{\epsilon,*}}
\le\frac{c_{\mathrm{par}}}{n}
\le\mathfrak R^Z_{n,d,\epsilon}.
\]
Together with the two large-sample cases, this establishes the lower
bound for every \(n,d\ge1\). Define
\[
C_\epsilon:=\max\{C_{\epsilon,\mathrm{up}},c_\epsilon\}.
\]
Then \(0<c_\epsilon\le C_\epsilon<\infty\). Since
\(r^{\mathrm{fix}}_{n,d}\ge0\), the upper bound already proved gives
\[
\mathfrak R^Z_{n,d,\epsilon}
\le
C_{\epsilon,\mathrm{up}}r^{\mathrm{fix}}_{n,d}
\le
C_\epsilon r^{\mathrm{fix}}_{n,d}.
\]
Thus, for every \(n,d\ge1\),
\[
c_\epsilon
\min\left\{
1,\frac1n+\left[\frac d{n\log(e+n)}\right]^2
\right\}
\le
\mathfrak R^Z_{n,d,\epsilon}
\le
C_\epsilon
\min\left\{
1,\frac1n+\left[\frac d{n\log(e+n)}\right]^2
\right\},
\]
as claimed.
% lean: CausalSmith.Stat.MarRareqLogfrontier.zeng_fixed_positivity_polynomial_frontier
\end{enumerate}
\end{proof}
